\documentclass[
  paper=a4,
  fontsize=12pt,
  DIV=16,
  headheight=52pt,
  footheight=45pt,
  headinclude,
  parskip=full,
]{scrartcl}

\usepackage{fontspec}
\setsansfont{Source Sans Pro}
\renewcommand{\familydefault}{\sfdefault}

\usepackage{polyglossia}
\setmainlanguage{german}

\usepackage{csquotes}

\usepackage{microtype}
\usepackage{graphicx}

\usepackage{scrlayer-scrpage}
\pagestyle{scrheadings}
\setkomafont{pagehead}{\normalfont}
\setkomafont{pagefoot}{\normalfont\footnotesize}
\ohead{\includegraphics[height=1.5cm]{pep.pdf}}
\ihead{%
  \large\bfseries PeP et al.\ e.\,V.\\%
  Begabtenförderungsprogramm –\\
  Going Abroad - Purdue%
}
\cfoot{}
\ifoot{%
  PeP et al.~e.\,V.\\
  Begabtenförderung\\
  \url{www.pep-dortmund.org}
}


\usepackage{titling}
\usepackage{booktabs}

\usepackage{xcolor}
\usepackage{calc}
\usepackage[colorlinks=true,urlcolor=blue!50!black]{hyperref}

\date{31. Juli 2014}

\begin{document}
\textbf{\Huge\sffamily Bachelorarbeit in Purdue, Indiana}\\[0.5\baselineskip]
\textbf{\Large\sffamily Eine Anleitung was alles im Vorfeld tun ist}

Du hast die Zusage zur Teilnahme am Going Abroad Programm für Purdue erhalten.
Mit der alleinigen Zusage ist es leider nicht getan.
Es gibt noch einige wichige Aufgaben, die du vorher erledigen musst.
Im folgenden werden dir die wichtigsten Fragen beantwortet.
Solltest du nicht ausreichend informiert sein, oder noch weitere Fragen haben,
kannst du dich an Henning Moldenhauer (CP-O1-190) persönlich oder per Mail
(henning.moldenhauer@pep-dortmund.org) wenden.

\subsubsection*{Wie bekomme ich das richtige Visum für die USA?}
Das für dich richtige Visum ist das J1 Visum.
Der Prozess ist etwas komplexer und wird dir im folgenden dargelegt.

Zuerst musst du die \textbf{Einverständniserklärung der Fakultät} von deinem
betreuenden Professor \textbf{unterschreiben lassen} und bei Henning im Büro
abgeben.
Diese Erklärung ist sehr wichtig und wird schnellstmöglich von uns benötigt.
Erst mit ihr können wir den Visumsprozess an der Purdue University starten, sowie
deine Unterkunft vor Ort organisieren.

Ebenfalls entscheidend für die Beantragung deines Visums ist der
\textbf{Arbeitsvertrag}, den du mit der Einverständniserklärung unterschrieben
bei uns abgeben musst.
% Noch etwas besser begründen.

Sobald wir die Dokumente unterschrieben zurück bekommen haben, werden diese an
Andreas Jung in die USA weitergeleitet.
Er wird dann an der Purdue Universität dafür sorgen, dass du dort auch angenommen
wirst. Dazu muss wird er bei verschiedenen Stellen, wie dem Business Office,
dem Visa and Immigration Office und dem Office of International Students and
Scholars (ISS), alles notwendige in die Wege leiten.
% werden hier noch weitere Dokumente benötigt?

Mit den unterschriebenen Dokumenten musst du uns auch eine
\textbf{Kopie des Reisepasses} einreichen, damit auf amerikamischer Seite die
Dokumente für das Visum mit den richtigen Daten ausgefüllt werden können.

Während dieser Vorbereitungen wirst du ein paar Mails erhalten, die du
ebenfalls zügig beantworten solltest.
Unter anderem wirst du dazu aufgefordert werden die
\textbf{MyISS J Request E-forms} auszufüllen.
% was genau für Mails? von wem? was wird in etwa gefragt?

\subsubsection*{Warum muss ich mich Krankenversichern?}
Eine Sommerschule in den USA ist nur mit einem J-1 Visum möglich.
Der Visaprozess erfordert zwingend das Vorhandensein einer
\textbf{US tauglichen Krankenversicherung}.
In den Staaten reicht deine Versicherung nicht aus und du musst selbst eine
Zusatzversicherung abschließen.
Hier gibt es verschiedene Möglichkeiten:
\begin{itemize}
  \item Der DAAD bietet über die Continentale (sitzt in Dortmund) solche
  Zusatzversicherungen für etwa 160 € im Monat an.
  \item Es kann auch eine Krankenversicherung in den USA zugekauft werden, hier
  gibt es zum Beispiel die folgenden Varianten:
  \begin{itemize}
    \item \url{www.compassstudenthealthinsurance.com/compare_international_insurance_plans.php}
    Eine US Krankenversicherung immer mit sehr viel Auswahl verbunden.
    Wir empfehlen eine most-inclusive insurance, um bei tatsächlicher Krankheit
    den hohen Kosten des US Gesundheitssystems vorzubeugen.
    Compass Benchmark Plus in der entsprechenden Altersklasse liegt - mit
    Selbstbeteiligung - bei etwa 200 € pro Monat.
    \item Das ISS von Purdue empfiehlt diese zwei Optionen:
    \begin{itemize}
      \item \url{www.compassstudenthealthinsurance.com/}
      \item \url{www.isoa.org/}
    \end{itemize}
  \end{itemize}
\end{itemize}

\subsubsection*{Wie komme ich nach West Lafayette, Indiana?}
Mit dem Flugzeug nach Indianapolis, Indiana. Und von da aus weiter mit dem Bus
nach West Lafayette.

\subsubsection*{Wo schlafe ich?}

\end{document}
